\chapter{Division des nombres relatifs}

\begin{itemize}
\item Effectuer les divisions suivantes :
\begin{itemize}
\it $(+56):(+7)$.
\it $(-56):(-8)$.
\it $(+64):(-16)$.
\it $(-360):(+12)$.
\it $(+105):(-7)$. 
\it $(-286):(-13)$.
\it $(+221):(-17)$.
\it $(-286):(-11)$.
\it $(+96):(+8)$. 
\end{itemize} 
\item Effectuer les divisions suivantes : 
\begin{itemize}
\it $\displaystyle(+15):\left(-\frac23\right)$.
\it $\displaystyle(-7,5):\left(-\frac45\right)$.
\it $\displaystyle(-1):\left(+\frac34\right)$.
\it $\displaystyle\left(-\frac45\right):(+1,4)$.
\it $\displaystyle\left(+\frac34\right):\left(-\frac{15}8\right)$.
\it $\displaystyle\left(-\frac73\right):\left(+\frac53\right)$.
\it $\displaystyle(+1):\left(+\frac23\right)$.
\it $\displaystyle\left(-\frac{8}{25}\right):\left(-\frac{14}5\right)$.
\it $\displaystyle\left(+\frac13\right):\left(-\frac14\right)$.
\end{itemize}
\item Effectuer les divisions suivantes : 
\begin{itemize}
\it $\displaystyle\frac{-12+4-72}{4}$.
\it $\displaystyle\frac{+35-75+25}{5}$.
\it $\displaystyle\frac{-12+144-108}{-12}$.
\it $\displaystyle\frac{7-0,5+2}{-2}$. 
\it $\displaystyle\frac{9+13-15}{11}$.
\it $\displaystyle\frac{-11+15-7,8}5$. 
\end{itemize}
\item Simplifier les rapports suivants et trouver leur valeur : 
\begin{itemize}
\it $\displaystyle\frac{-63}{70}$.
\it $\displaystyle\frac{-308}{484}$.
\it $\displaystyle\frac{273}{468}$. 
\it $\displaystyle\frac{48}{-96}$.
\it $\displaystyle\frac{-169}{-26}$.
\it $\displaystyle\frac{-365}{-146}$.
\end{itemize}
\item Effectuer les opérations suivantes : 
\begin{itemize}
\it $\displaystyle\frac{-3+\frac34-\frac13}{5+\frac25-\frac23}$.
\it $\displaystyle\frac{7-\frac32+\frac38}{2-\frac52-\frac34}$.
\it $\displaystyle\frac{1-\frac57-\frac49}{1-\frac53-\frac7{15}}$.
\it $\displaystyle\frac{-3-0,5-0,75}{-\frac23-\frac34}$. 
\end{itemize}
\item Effectuer les opérations suivantes : 
\begin{itemize}
\it $\displaystyle\frac{-3+5-6}{2+7-1}\times\frac{4-3+1}{2-3-1}\times \frac{-5-7}{10}$.
\it $\displaystyle\frac{\frac12-\frac23}{\frac34-1}\times \frac{\frac56+\frac13}{1-\frac45}\times \frac{-\frac78+1}{\frac32}.$
\end{itemize}
\item Effectuer les opérations suivantes : 
\begin{itemize}
\it $\displaystyle\frac{0,5}{\frac35-2}+\frac{1}{4-\frac{11}7}-\frac3{-7+\frac13}$.
\it $\displaystyle 3- \frac{1-\frac12}{1+\frac12}$.
\it $\displaystyle \frac{\frac35}{\frac23-\frac45}-
\frac{\frac12}{\frac18-1}+\frac{\frac12}{\frac1{16}-2}$.
\it $\displaystyle\frac{1-\frac13+\frac1{1+\frac13}}{1-\frac13-\frac1{1+\frac13}}$
\end{itemize}
\item Le nombre $x$ prend toutes les valeurs entières de $-5$ à $+10$. Dresser le tableau de correspondance entre $x$ et $y$ dans les cas suivants : 
\begin{itemize}
\it $\displaystyle y=\frac{3x}5+1$.
\it $\displaystyle y=\frac{-2x}3+4$.
\it $\displaystyle y=\frac{-3x}2-4$.
\it $\displaystyle y=\frac{-7x}5-10$.
\it $\displaystyle y=\frac{2x+3}5$.
\it $\displaystyle y=\frac{-3x+1}4$.
\it $\displaystyle y=\frac{3x-7}2$.
\it $\displaystyle y=\frac{-5x-3}3$.
\end{itemize}
\end{itemize}